\chapter{Integrace spisové služby}
V původní analýze informačního systému\cite{analyza-PMSIS} školy byla spisová služba uvedena pouze okrajově a byl popsán pouze účel, avšak bez hlubšího zamyšlení nad její prioritou a bez podrobného rozboru toho, co integrace spisové služby plně zahrnuje. Tento přístup odrážel původní předpoklad, že zatížení spojené s implementací tohoto komplexního systému by pro školu nepředstagovalo největší prioritu.

Avšak s přibývajícími zkušenostmi a dalším zkoumáním administrativních procesů ve škole došlo k přehodnocení této situace. Bylo zjištěno, že administrativní zátěž na zaměstnance, včetně třídních učitelů, je výrazně vyšší, než se původně předpoklávalo. Toto je způsobeno rozsáhlými byrokratickými požadavky na zpracování a správu dokumentů, které by mohla integrace spisové služby výrazně zjednodušit prostřednictvím automatizace generování dokumentů.

Pro přiblížení je zde příklad procesu vyřízení pochvaly třídního učitele nezletilého žáka.
\begin{figure}[H]
    \centering
    \includegraphics[width=\textwidth/2]{ac-priklad-pochvala-tridniho-ucitele.png}
    \caption{Příklad procesu vydání pochvaly třidního učitele}
    \label{fig:priklad-pochvala-tridniho-ucitele}
\end{figure}

Tento příklad je jeden z celé škály. Zde uváním všechny dokumenty z kategorie výchovných opatřední, které musí spravovat řadový třídní učitel:
\begin{itemize}
    \item zkouška pro doplnění klasifikace - oznámení;
    \item pochvala třídního učitele;
    \item důtka třídního učitele;
    \item napomenutí třídního učitele;
    \item upozornění na velkou absenci a na možnost vyžadování omlouvání absence pouze 
    lékařem nebo jiným úředním způsobem;
    \item omlouvání avsence pouze lékařem nebo jiným úředním způsobem;
    \item upozornění na zaostávání;
    \item upozornění na neúčast na cvičeních a neodevzdávání zpráv ze cvičení;
\end{itemize}

Některé z těchto dokumentů mají další specifické podverze. Další typy a kategorie dokumentů jsou pro například pro ředitele školy. V uvedené kategorii celou manipulaci vyřizuje jeho zástupce pro třídu - tedy třídní učitel.

V současné době se jednací čísla generují manuálně prostřednictvím administrativního pracovníka, což je proces, který je časově i personálně náročný a náchylný k lidským chybám. Tento tradiční přístup zahrnuje několik kroků, od žádosti o číslo, přes jeho zápis do fyzických nebo digitálních dokumentů, až po archivaci. Výsledkem je, že celý proces vyřízení jednoduchých úkonů, jako je vydání pochvaly třídního učitele, trvá déle, než je žádoucí, a to hlavně kvůli administrativní zátěži, která s tím souvisí.

Existuje možnost plné integrace spisové služby do ŠIS, což by výrazně zjednodušilo a urychlilo proces generování jednacích čísel a správu dokumentů. Pro úspěšnou integraci je však nezbytné, aby systém splňoval požadavky Národního standardu pro elektronické systémy spisové služby (NSESSS)\cite{nsesss}, což je klíčový dokument, který stanovuje pravidla pro správu, ukládání, archivaci a ochranu elektronických dokumentů ve veřejné správě.

NSESSS je komplexní sada specifikací, která má za cíl zvýšit efektivitu a bezpečnost v elektronickém zpracování dokumentů. Standard pokrývá širokou škálu témat, od formátů dokumentů, přes jejich klasifikaci, až po procesy schvalování a ukládání. Díky tomu může spisová služba, která je s NSESSS kompatibilní, značně zjednodušit běznou administrativu a zároveň zajistit vysokou úroveň ochrany a dostupnosti dat.

V kontextu implementace způsobů získávání spisových čísel se nabízí dvě základní cesty, jak dosáhnout větší efektivity a zjednodušení. Tyto cesty reflektují rozdílné přístupy k řešení výzvy spojené s integrací a získáváním spisových čísel, přičemž každá z nich má své specifické výhody i omezení.

\section{Implementace NSESSS}
První možností je implementace NSESSS, což je komplexní a normativně náročný proces, který by umožnil plně automatizované získávání spisových čísel ze spisové služby. Tento přístup vyžaduje hluboké porozumění a studium XML struktur, které jsou základem pro správnou implementaci spisové služby v souladu s NSESSS. Avšak právě rozsáhlé požadavky na technické a legislativní znalosti činí tento přístup časově a kapacitně náročným. Ve světle omezených financí a potřeby zavedení systému, který bude efektivní z hlediska ceny a výkonu, se tento způsob jeví jako méně vhodný.

Příklad fungování systému s implementací NSESSS
\begin{figure}[H]
    \centering
    \includegraphics[width=\textwidth-60pt]{ac-nsesss-pochvala.png}
    \caption{Příklad procesu vydání pochvaly třidního učitele pomocí NSESSS. Modře jsou vyznačeny aktivity zpracovávané člověkem a zeleně softwarovou část.}
    \label{fig:ac-nsesss-pochvala}
\end{figure}

\section{Implementace vícestupňového formuláře}
Druhou cestou je vytvoření vícestupňového formuláře, který představuje více praktické a flexibilní řešení. Tento systém je navržen tak, aby automaticky notifikoval příslušného zaměstnance v kanceláři a poskytl mu přehledné informace potřebné pro rychlé zpracování žádosti o spisové číslo. Tento přístup minimalizuje potřebu rozsáhlého studia NSESSS pro programátora a současně zlepšuje efektivitu zpracování vydávání dokumentů.

Dle kontaktu s ředitelem školy jsem provedl analýzu uživatelských požadavků a na jeho základě jsem identifikoval dvě možná řešení pro implementaci tohoto vícestupňového formuláře, které jsou nyní předmětem zájmu. Tento proces je klíčový pro efektivní uživatelsky-centristický nárh systému (UCD).

\subsection{Dvoufázový formulář}
Dvoufázový formulář je základním řešením pro automatizaci procesu žádosti o spisové číslo. Jeho hlavní předností je jednoduchost a přímá aplikace.
\begin{enumerate}
    \item V první fázi je žádost o vydání dokumentu žadatelem, pro který bude získáváno spisové číslo, prostřednictvím elektronického formuláře, který žadatel vyplní dle specifických požadavků pro daný dokument. Po odeslání formuláře systém automaticky notifikuje zodpovědného zaměstnance v kanceláři, který přidělí spisové číslo.

    \item Ve druhé fázi si zodpovědný zaměstnanec zobrazí data z odeslaného formuláře a zadá je do systému spisové služby. Ta přidělí dokumentu spisové číslo zaměstnanec manuálně zapíše jednací číslo do formuláře. Systém notifikuje žadatele o zpracování a může si hotový dokument zobrazit, vytisknout a dále s ním manipulovat dle potřeb.
\end{enumerate}

\begin{figure}[H]
    \centering
    \includegraphics[width=\textwidth/2]{ac-dvoufaz-pochvala.png}
    \caption{Příklad procesu vydání pochvaly třidního učitele pomocí dvoufázového formuláře. Modře jsou vyznačeny aktivity zpracovávané člověkem a zeleně softwarovou část.}
    \label{fig:ac-dvoufaz-pochvala}
\end{figure}

\subsection{Třífázový formulář}
Třífázový formulář představuje inovativní rozšíření dvoufázového modelu tím, že přidává specifický prvek interakce do stávajícího procesu - hlasování o některých druzích důtek a pochval a dalších dokumentech pedagogickým sborem. Ředitel považuje za důležité hlasovat o ředitelských pochvalách, důtkách, povoleních opakovat ročník, podmínečných vyloučení a dalších výchovných opatřední celý sbor a svěřuje tímto tuto svou pravomoc přímo pedagogickým pracovníkům. Žáky bude moct třídní učitel nominovat na uvedené na pochvaly, důtky a další výchová opatřední schvalované ředitelem prostřednictvím IS.

\begin{enumerate}
    \item V tomto modelu první fáze zůstává identická s první fází dvoufázového formuláře, což znamená, že začíná standardním formulářem, který připraví žádost o spisové číslo.

    \item Druhá fáze, která je jedinečná pro třífázový model, přináší nový rozměr – platformu pro hlasování. Tato platforma umožňuje zapisovateli porad urychleně zapsat počty hlasujících PRO, PROTI, ZDRŽEL SE do IS a ten rozhodne na základě převahy skupin PRO a PROTI o přechodu do třetí fáze, či skartace záznamu.

    \item Třetí a závěrečná fáze opět navazuje na standardní proces dvoufázového formuláře, kde po hlasování následuje automatická notifikace zodpovědného zaměstnance a zaznamenání spisového čísla do systému.
\end{enumerate}
Tímto způsobem se elektronicky zaznamená hlasování o rozhodnutí, které se nyní dělá ručně v papírové formě a zároveň systém ihned  

\begin{figure}[H]
    \centering
    \includegraphics[width=\textwidth-120pt]{ac-trojfaz-pochvala.png}
    \caption{Příklad procesu vydání pochvaly třidního učitele pomocí třífázového formuláře. Modře jsou vyznačeny aktivity zpracovávané člověkem a zeleně softwarovou část.}
    \label{fig:ac-trojfaz-pochvala}
\end{figure}